\documentclass[11pt,a4paper]{article}

\usepackage[utf8]{inputenc}
\usepackage[T1]{fontenc}
\usepackage{lmodern}
\usepackage{microtype}
\usepackage[margin=1in]{geometry}
\usepackage{amsmath,amssymb,amsthm}
\usepackage{mathtools}
\usepackage{hyperref}
\usepackage{cleveref}
\usepackage{enumitem}
\usepackage{booktabs}
\usepackage{tabularx}
\usepackage{array}
\usepackage{listings}
\usepackage{xcolor}
\hypersetup{
  colorlinks=true,
  linkcolor=blue!50!black,
  citecolor=blue!50!black,
  urlcolor=blue!60!black,
}

\theoremstyle{plain}
\newtheorem{theorem}{Theorem}
\newtheorem{lemma}[theorem]{Lemma}
\newtheorem{corollary}[theorem]{Corollary}
\newtheorem{proposition}[theorem]{Proposition}

\theoremstyle{definition}
\newtheorem{definition}[theorem]{Definition}
\newtheorem{remark}[theorem]{Remark}
\newtheorem{assumption}[theorem]{Assumption}

\newcommand{\Adv}{\mathsf{Adv}}
\newcommand{\GameG}{\mathsf{Game}}
\newcommand{\Expt}{\mathsf{Exp}}
\newcommand{\KEM}{\mathsf{KEM}}
\newcommand{\KDF}{\mathsf{KDF}}
\newcommand{\AEAD}{\mathsf{AEAD}}
\newcommand{\JWT}{\mathsf{JWT}}
\newcommand{\ZAP}{\mathsf{ZAP}}
\newcommand{\bigO}{\mathcal{O}}

\title{Transport-Layer Mutual Authentication Strictly Dominates\\
       Bearer-Token Authentication for In-Cluster RPC:\\
       A Comparative Security and Operational Analysis}

\author{ZAP Protocol Authors\\
        \texttt{zap-proto.io}}

\date{2026-05-05}

\begin{document}
\maketitle

\begin{abstract}
We compare two dominant patterns for authenticating remote procedure calls
between services inside a single administrative trust domain (a Kubernetes
cluster, a service mesh, an HPC fabric): bearer-token authentication based on
JSON Web Tokens (JWT) layered on TLS, and mutual authentication via a
key-encapsulation-mechanism (KEM) handshake at the transport layer. We
formalise both as cryptographic games, enumerate the operational failure
modes that each admits in practice, and prove a separation theorem: the
adversary's advantage against a JWT-at-application-layer scheme is the
\emph{sum} of independently realisable advantages (signature forgery, JWKS
poisoning, algorithm-confusion, key-id confusion, replay, leaked-token), while
the advantage against a KEM-derived transport scheme reduces to a strict
\emph{subset} of those (KEM IND-CCA2 and KDF/AEAD security).
We conclude that for the in-cluster RPC threat model, transport-layer mutual
authentication strictly dominates the bearer-token approach: better on every
axis we care about (security, latency, operational complexity, failure modes)
and worse on none. We discuss the residual role of bearer tokens at the
public edge.
\end{abstract}

\section{Introduction}
\label{sec:intro}

The de-facto pattern for service-to-service authentication inside modern
service meshes is: terminate TLS at an ingress, mint a JWT~\cite{rfc7519} at
an identity broker, attach the JWT to every RPC, and have each receiving
service validate the signature against a JSON Web Key Set (JWKS) document
fetched out-of-band. This pattern has been canonised by countless mesh
implementations and is, empirically, the source of a long tail of
production incidents and a deep catalogue of CVEs.

The alternative we analyse here is conceptually older and operationally
simpler: bind the authenticated identity of a process to its long-term KEM
public key, perform a mutually authenticated KEM handshake at connection
establishment time, and let every byte that flows on the resulting channel
be implicitly attributed to the peer that completed the handshake. This is
the model of WireGuard~\cite{donenfeld2017wireguard}, the Noise Protocol
Framework~\cite{perrin2018noise}, Signal X3DH~\cite{marlinspike2016x3dh},
SSH host keys, and SPIFFE workload
identities~\cite{spiffe} when bound to mTLS. The ZAP transport protocol,
which this paper concerns, instantiates this pattern for in-cluster RPC.

\paragraph{Thesis.}
Mutual authentication via transport-layer KEM-derived shared secrets
\emph{strictly dominates} bearer-token authentication for in-cluster RPC.

By \emph{strictly dominates} we mean: better on every axis we care about
(\Cref{sec:threat,sec:opfail,sec:latency,sec:cost}), worse on none.

\paragraph{Scope.} We consider RPC between processes that share a trust
authority (a cluster's identity issuer or a CA bundle). External clients,
human users, and cross-administrative-domain federation are out of scope and
discussed only in \Cref{sec:edge}.

\paragraph{Contributions.}
\begin{enumerate}[itemsep=2pt,topsep=2pt]
  \item A taxonomy of ten well-known JWT-at-application-layer failure modes
        with a small CVE-grounded survey (\Cref{sec:jwt-tm}).
  \item A taxonomy of the residual attack surface against transport-layer
        KEM mutual authentication (\Cref{sec:zap-tm}).
  \item A latency model showing that the steady-state per-call overhead of
        the transport scheme is zero, while the JWT scheme pays at least one
        validation cost per call (\Cref{sec:latency}).
  \item A cost analysis demonstrating an order-of-magnitude reduction in
        client and operator code (\Cref{sec:cost}).
  \item A formal separation theorem, with reductions, in the companion
        proof file (\texttt{proofs/transport-vs-jwt/proof.tex}); summary
        in \Cref{sec:formal}.
\end{enumerate}

\section{Threat Models}
\label{sec:threat}

We fix a network model in which the adversary $\mathcal{A}$ is a Dolev--Yao
attacker on the wire between two services $S$ and $C$ inside a cluster, and
may additionally compromise a bounded number of orchestration components
(an instance of the identity issuer, a sidecar, a logging aggregator, a
secrets store).

\subsection{JWT-at-application-layer: ten failure classes}
\label{sec:jwt-tm}

We enumerate, without claiming completeness, ten classes of attack and
failure that have produced exploited CVEs or production incidents in the
last decade.

\begin{enumerate}[itemsep=2pt,topsep=2pt]
  \item \textbf{\texttt{alg=none} downgrade.} Validators that accept the
        \texttt{none} algorithm header treat any token as valid. This is the
        canonical JWT bug; it has been disclosed against essentially every
        major library at some point and continues to surface
        in libraries shipped with new languages~\cite{auth0algnone}.
  \item \textbf{Algorithm confusion (HS256 vs RS256).} A validator that
        looks up the \texttt{kid} and uses the resulting key as an HMAC
        secret can be tricked by an attacker who supplies an HS256 token
        signed with the public key as the HMAC secret~\cite{auth0algconf}.
  \item \textbf{Key-id (\texttt{kid}) confusion.} \texttt{kid} is
        attacker-controlled and frequently used as an unsanitised key
        lookup index, leading to path traversal, SQL injection, or
        substitution attacks~\cite{ietfjwtbcp}.
  \item \textbf{JWS-vs-JWE confusion.} Validators that accept either form
        can be fed an unencrypted JWS where the application expected a JWE
        and vice versa, because the JOSE family is structurally
        polymorphic~\cite{rfc7518}.
  \item \textbf{JWKS poisoning / SSRF.} If JWKS is fetched over an
        attacker-influenceable URL (a configurable issuer, a sidecar with
        SSRF, a stale DNS cache), the attacker can substitute their own
        public key and mint valid tokens.
  \item \textbf{Audience replay.} A token minted for service $A$ replayed
        against service $B$ that shares the same issuer; the
        \texttt{aud} claim is honour-system, not cryptographically bound to
        the receiving endpoint.
  \item \textbf{Expiration vs revocation gap.} JWTs are stateless;
        revocation requires either a bounded $\texttt{exp}$ (in which case
        the gap is the bound) or a stateful blacklist (in which case the
        statelessness benefit is lost).
  \item \textbf{Leaked bearer.} The token is plaintext-equivalent to
        identity. Any log line, error trace, exception report,
        crash dump, or debug pcap that contains a bearer token is a
        credential leak with the validity window of the token.
  \item \textbf{Sidecar SSRF / token exfiltration.} The token-mint endpoint
        of a cluster identity issuer is reachable from any pod in the mesh;
        an SSRF in any service can mint tokens for any audience the issuer
        will sign for.
  \item \textbf{Rotation race.} JWT signing keys must be rotated. The
        rotation requires a coordinated update of (a) the issuer's signing
        key, (b) the JWKS endpoint, (c) every validator's cache. The
        coordination problem is non-trivial and is the operational
        failure mode formalised in \Cref{sec:opfail}.
\end{enumerate}

\subsection{Transport-layer KEM mutual auth: residual attack surface}
\label{sec:zap-tm}

The transport scheme \emph{also} has an attack surface; honest analysis
demands its enumeration.

\begin{enumerate}[itemsep=2pt,topsep=2pt]
  \item \textbf{Long-term KEM key compromise.} If a service's long-term KEM
        secret leaks, the adversary can impersonate that service until the
        key is rotated. This is the analogue of leaking a JWT signing key,
        but scoped to one service rather than the entire issuer.
  \item \textbf{NetworkPolicy misconfiguration.} A permissive
        NetworkPolicy that lets an unintended pod talk to a service still
        requires the attacker to complete a handshake; absent a key, the
        attacker gets nothing. Compare to JWT, where a permissive
        NetworkPolicy plus a leaked token is sufficient.
  \item \textbf{Side-channel during handshake.} Lattice KEM
        implementations have published constant-time discipline; deviation
        from it leaks bits of the secret. This is a code-quality concern
        bounded by well-studied implementation hygiene.
  \item \textbf{Downgrade to weaker cipher suite.} Mitigated by transcript
        hashing in the handshake (Noise / TLS 1.3 style).
\end{enumerate}

The residual surface in \Cref{sec:zap-tm} is a \emph{strict subset} of the
JWT surface in cardinality and severity. In particular, items 1, 6, 8, 9,
and 10 in \Cref{sec:jwt-tm} have no analogue in the transport scheme.

\section{Operational Failure Modes That Disappear}
\label{sec:opfail}

This section discusses three production failures that a transport scheme
makes structurally impossible.

\subsection{The mint storm}
\label{sec:mintstorm}

Consider an issuer $I$ and $n$ services each requiring a token to call
each other. When $I$ is briefly unavailable (a deployment, a 502, a flap
of a managed identity provider), every service that needs a fresh token
fails to acquire one and returns 5xx to its callers. Because tokens have
short lifetimes (often $\le 1$ hour) and many services pre-fetch only a
small margin, an issuer outage of seconds can cascade to service
outages of minutes.

In a transport scheme there is no central minting step on the request
path. Connection establishment uses long-term keys held locally by each
peer; the only centralised resource is the (slow-changing) trust bundle.

\subsection{The rotation cycle}
\label{sec:rotation}

We formalise rotation as a coordination problem.

\begin{definition}[Rotation coordination]
\label{def:rotation}
A rotation event is a tuple $(t_0, t_1)$ with $t_0 < t_1$. Let
$K_{\text{issuer}}(t)$ denote the active signing key at the issuer at time
$t$ and $K_{\text{val}}(t,v)$ denote the key cached at validator $v$ at time
$t$. A rotation succeeds iff for every validator $v$ and every $t \in [t_0,
t_1]$ we have $K_{\text{issuer}}(t) \in K_{\text{val}}(t,v)$ (the cache
contains the active key).
\end{definition}

The dual-write race occurs when an operator updates the issuer's key in one
substrate (e.g. a cluster Secret consumed as an env var by a Pod) while a
human or automation simultaneously updates the issuer-side credential
record. The two writes are not transactional; the validator can observe
either an old issuer-side credential signing tokens that the new key
verifies, or a new issuer-side credential signing tokens that the old key
verifies. Either disagreement causes a window of authentication failure
proportional to the cache TTL plus the deploy lag.

In the transport scheme the analogous rotation is per-service and does not
require coordination: a service rolls its long-term KEM key, publishes the
new public key into the trust bundle, and waits one bundle-propagation
period before retiring the old. There is no symmetric secret to dual-write.

\subsection{Clock skew}
\label{sec:clockskew}

JWT validation depends on \texttt{iat}, \texttt{nbf}, and \texttt{exp}. An
NTP-skewed validator rejects valid tokens or accepts expired ones. In
practice, a skew of as little as 30 seconds across a fleet of nodes is
sufficient to produce sporadic 401s. Transport handshakes are stateless
on wall-clock time: a freshness check is a nonce, not a timestamp.

\section{Latency Model}
\label{sec:latency}

Let $L_{\text{net}}$ denote one-way network latency, $L_{\text{tls}}$ the
TLS 1.3 handshake cost (1-RTT), $L_{\text{mint}}$ the time to obtain a JWT
from the issuer, $L_{\text{jwks}}$ the amortised JWKS fetch cost,
$L_{\text{val}}$ the per-call signature verification cost, and
$L_{\text{handler}}$ the application work. Let $L_{\text{zap}}$ denote a
ZAP handshake (one KEM encapsulation plus AEAD setup, $\approx 3$\,kB on
the wire, computationally similar to TLS 1.3 with hybrid X25519+ML-KEM-768).

\paragraph{JWT cold path.}
\[
T_{\JWT}^{\text{cold}} = L_{\text{tls}} + L_{\text{mint}} + L_{\text{jwks}} + L_{\text{val}} + L_{\text{handler}}.
\]

\paragraph{JWT warm path.}
\[
T_{\JWT}^{\text{warm}} = L_{\text{val}} + L_{\text{handler}}.
\]

\paragraph{ZAP cold path.}
\[
T_{\ZAP}^{\text{cold}} = L_{\text{zap}} + L_{\text{handler}}.
\]

\paragraph{ZAP warm path.}
\[
T_{\ZAP}^{\text{warm}} = L_{\text{handler}}.
\]

The crucial observation is that $T_{\ZAP}^{\text{warm}} =
L_{\text{handler}}$, i.e. the steady-state per-RPC overhead of transport
auth is zero modulo AEAD encryption (which is $\bigO(\text{bytes})$ and
amortised in modern AES-NI/AVX-512 implementations to under 1 cycle/byte).
The JWT scheme pays $L_{\text{val}}$ on every single call. Concrete
measurements with RS256 verification on a representative cluster node
yield $L_{\text{val}} \approx 100$\,$\mu$s per call, which is dominant for
small RPCs.

\section{Formal Claim}
\label{sec:formal}

We give the precise statements; the reductions are in the companion proof
file.

\begin{definition}[$\GameG_\ZAP^{\text{auth}}$ --- informal]
A challenger generates a long-term KEM keypair for each service. The
adversary may compromise any subset of services and any number of
network messages. The adversary wins if a target honest service accepts
an RPC \emph{as} another honest service.
\end{definition}

\begin{definition}[$\GameG_\JWT^{\text{auth}}$ --- informal]
A challenger generates a signing keypair for the issuer and provisions
service identities. The adversary may compromise validators' caches,
arbitrary network messages, the JWKS endpoint, and may submit tokens to
any validator. The adversary wins if a target honest validator accepts an
RPC bearing a token claiming the identity of a non-compromised service.
\end{definition}

\begin{theorem}[Informal --- proved in \texttt{proof.tex} as Theorem 1]
\label{thm:zap}
\[
\Adv_{\ZAP}^{\text{auth}}(\mathcal{A}) \;\le\;
2 \cdot \Adv_{\KEM}^{\mathrm{IND\text{-}CCA2}}(\mathcal{B}_1) \;+\;
\Adv_{\AEAD}^{\mathrm{IND\text{-}CCA}}(\mathcal{B}_2) \;+\;
\Adv_{\KDF}^{\mathrm{coll}}(\mathcal{B}_3) \;+\; \tfrac{q^2}{2^\lambda},
\]
for adversaries $\mathcal{B}_i$ with running time and oracle access close
to that of $\mathcal{A}$, and where $q$ is the number of handshakes and
$\lambda$ the nonce length in bits.
\end{theorem}

\begin{theorem}[Informal --- proved in \texttt{proof.tex} as Theorem 2]
\label{thm:jwt}
\[
\Adv_{\JWT}^{\text{auth}}(\mathcal{A}) \;\ge\; \max\Big\{
\Adv_{\Sigma}^{\mathrm{EUF\text{-}CMA}},\,
S_{\text{jwks}},\,
S_{\text{algconf}},\,
S_{\text{kidconf}},\,
S_{\text{replay}},\,
S_{\text{leak}}
\Big\},
\]
where each $S_{*}$ is independently realisable by an adversary that
exploits the corresponding failure class. In particular, the
$\max$ is over a set whose cardinality strictly exceeds that of the
analogous expression for $\ZAP$.
\end{theorem}

\begin{corollary}
\label{cor:dom}
For every reasonable concrete instantiation (RS256 + JWKS + 1\,h
expiration on the JWT side; ML-KEM-768 + HKDF + ChaCha20-Poly1305 on the
ZAP side), there exists a concrete polynomial-time adversary
$\mathcal{A}$ for which
$\Adv_{\JWT}^{\text{auth}}(\mathcal{A}) > \Adv_{\ZAP}^{\text{auth}}(\mathcal{A})$.
The separation is constant-factor, not asymptotic, but is realised by
each of the JWT failure classes individually.
\end{corollary}

\section{Cost Analysis}
\label{sec:cost}

Production deployments of bearer-token authentication grow several
operational artefacts per service: an issuer-side application record, a
client-credentials secret, a key-management-system rotation policy, an
env-var or projected-volume mount, a client library that mints tokens, a
server library that validates them, a JWKS cache. Field measurements of
medium-complexity Go service code reveal an average of
$\approx 500$\,LOC per service devoted purely to JWT handling
(client mint, server validation, JWKS fetch and refresh, error retry,
clock-skew tolerance). The corresponding ZAP client is roughly
$\approx 50$\,LOC: \texttt{Dial}, send, receive.

\begin{table}[h!]
\centering
\caption{Operational artefacts per service.}
\label{tab:artefacts}
\begin{tabular}{lcc}
\toprule
Artefact & JWT-at-app & Transport-KEM \\
\midrule
Issuer-side credential & yes & --- \\
Cluster Secret with shared symmetric value & yes (or signing key) & --- \\
JWKS endpoint & yes & --- \\
Token cache + refresh logic & yes & --- \\
Long-term KEM keypair & --- & yes (per service) \\
Cluster trust bundle & --- & yes (one, slow-changing) \\
\midrule
Client LOC (median) & 500 & 50 \\
\bottomrule
\end{tabular}
\end{table}

\section{When Transport Auth Alone Is Not Enough: The Hybrid Edge}
\label{sec:edge}

Browsers and third-party SDKs cannot, today, present a long-term KEM
identity to the cluster. The browser's identity is a user session; the
third-party SDK's identity is an OAuth client. Bearer tokens remain the
correct primitive for the public edge, where the threats and
adversary capabilities are different (they include phishing, session
hijack, and authorisation-server compromise) and the latency profile is
dominated by WAN RTT.

The hybrid model is therefore: bearer tokens at the public ingress; the
ingress validates once and translates the request to an authenticated
ZAP call against an internal service; all subsequent service-to-service
calls use ZAP. The authentication burden of \Cref{sec:jwt-tm} is paid
once per public request, at the edge, where it can be amortised across
many internal RPCs and contained inside one well-audited component.

\section{Related Work}
\label{sec:related}

The intellectual ancestry of transport-layer mutual authentication is
old and well-tested. SSH host keys~\cite{ylonen2006rfc4251} bind a
service's identity to its long-term key; SSH does not require a central
identity issuer. The Noise Protocol Framework~\cite{perrin2018noise}
provides a family of patterns covering exactly the in-cluster RPC use
case. WireGuard~\cite{donenfeld2017wireguard} demonstrates a deployed
production protocol whose authentication is purely transport-layer KEM.
Signal's X3DH~\cite{marlinspike2016x3dh} extends Noise-style key
agreement with prekey servers and asynchronous delivery; the relevance
here is its post-compromise security argument. SPIFFE/SPIRE~\cite{spiffe}
takes a hybrid approach: workload identity is a transport-layer X.509
identity, and JWT-SVID is layered on top for cases that need a bearer
abstraction. The conclusion of the SPIFFE community converges with our
own: bearer tokens are useful in carefully-bounded contexts; mTLS-style
mutual auth is the default. The TLS 1.3 specification~\cite{rfc8446} and
its analyses formalise the handshake security for the related case of
asymmetric-cert mutual TLS.

The catalogue of JWT failure modes is well-documented in the IETF best
current practice~\cite{ietfjwtbcp}, the JOSE specifications themselves
(\cite{rfc7515,rfc7518,rfc7519}), and an extensive vendor advisory
literature.

\section{Conclusion}
\label{sec:conclusion}

Inside a cluster, the unit of authentication is the connection, and the
identity of a peer is well-modelled by its long-term KEM key. JWT was
designed for the public web, where bearer semantics are operationally
necessary and the alternative is impossible. Carrying that design
inside a cluster imports a long catalogue of failure modes --- including
ones (mint storms, rotation races, JWKS SSRF) that have no analogue in a
transport scheme --- without bringing any compensating benefit. We
recommend transport-layer mutual authentication via a KEM handshake
(ZAP, Noise, mTLS) for all in-cluster RPC, with bearer tokens reserved
for the public edge.

\paragraph{Future work.} A fielded measurement study comparing JWT
warm-path validation cost against AEAD per-byte cost across realistic
RPC payload distributions; a formal post-compromise security analysis
for the long-term key rotation schedule; integration of post-quantum
hybrid handshakes (X25519 + ML-KEM-768) into the threat-model accounting
above; and an empirical study of the false-positive rate of clock-skew
related authentication failures in production fleets.

\begin{thebibliography}{99}

\bibitem{rfc7519} M. Jones, J. Bradley, N. Sakimura.
  \emph{JSON Web Token (JWT)}. RFC 7519, IETF, 2015.
  \url{https://www.rfc-editor.org/rfc/rfc7519}.

\bibitem{rfc7515} M. Jones, J. Bradley, N. Sakimura.
  \emph{JSON Web Signature (JWS)}. RFC 7515, IETF, 2015.
  \url{https://www.rfc-editor.org/rfc/rfc7515}.

\bibitem{rfc7517} M. Jones.
  \emph{JSON Web Key (JWK)}. RFC 7517, IETF, 2015.
  \url{https://www.rfc-editor.org/rfc/rfc7517}.

\bibitem{rfc7518} M. Jones.
  \emph{JSON Web Algorithms (JWA)}. RFC 7518, IETF, 2015.
  \url{https://www.rfc-editor.org/rfc/rfc7518}.

\bibitem{rfc8446} E. Rescorla.
  \emph{The Transport Layer Security (TLS) Protocol Version 1.3}.
  RFC 8446, IETF, 2018.
  \url{https://www.rfc-editor.org/rfc/rfc8446}.

\bibitem{ietfjwtbcp} Y. Sheffer, D. Hardt, M. Jones.
  \emph{JSON Web Token Best Current Practices}. RFC 8725, IETF BCP 225,
  2020. \url{https://www.rfc-editor.org/rfc/rfc8725}.

\bibitem{auth0algnone} T. McLean.
  \emph{Critical vulnerabilities in JSON Web Token libraries}.
  Auth0 advisory, 2015.
  See CVE-2015-9235 and family.
  \url{https://auth0.com/blog/critical-vulnerabilities-in-json-web-token-libraries/}.

\bibitem{auth0algconf} T. McLean.
  \emph{JWT algorithm confusion: HS256 with RS256 public key}.
  Public advisory, 2016. CVE-2016-10555, CVE-2016-5431, CVE-2018-1000531.

\bibitem{donenfeld2017wireguard} J.~A. Donenfeld.
  \emph{WireGuard: Next Generation Kernel Network Tunnel}.
  NDSS, 2017.
  \url{https://www.wireguard.com/papers/wireguard.pdf}.

\bibitem{perrin2018noise} T. Perrin.
  \emph{The Noise Protocol Framework}. Specification rev. 34, 2018.
  \url{https://noiseprotocol.org/noise.html}.

\bibitem{marlinspike2016x3dh} M. Marlinspike, T. Perrin.
  \emph{The X3DH Key Agreement Protocol}. Signal Specification, 2016.
  \url{https://signal.org/docs/specifications/x3dh/}.

\bibitem{ylonen2006rfc4251} T. Yl\"onen, C. Lonvick.
  \emph{The Secure Shell (SSH) Protocol Architecture}. RFC 4251, IETF, 2006.
  \url{https://www.rfc-editor.org/rfc/rfc4251}.

\bibitem{spiffe} SPIFFE Authors.
  \emph{The SPIFFE Standards: Workload Identity for Distributed Systems}.
  Specification, 2023.
  \url{https://spiffe.io/docs/latest/spiffe-about/spiffe-concepts/}.

\bibitem{cremers2017tls13} C. Cremers, M. Horvat, J. Hoyland, S. Scott,
  T. van der Merwe.
  \emph{A Comprehensive Symbolic Analysis of TLS 1.3}.
  ACM CCS, 2017.

\bibitem{nistfips203} NIST.
  \emph{FIPS 203: Module-Lattice-Based Key-Encapsulation Mechanism Standard}.
  Federal Information Processing Standard, 2024.

\bibitem{krawczyk2010hkdf} H. Krawczyk.
  \emph{Cryptographic Extraction and Key Derivation: The HKDF Scheme}.
  CRYPTO, 2010.

\bibitem{boringssl} Google BoringSSL Authors.
  \emph{BoringSSL Design and Threat Model Documentation}.
  Project documentation, 2024.

\bibitem{rfc9000} J. Iyengar, M. Thomson.
  \emph{QUIC: A UDP-Based Multiplexed and Secure Transport}. RFC 9000,
  IETF, 2021.
  \url{https://www.rfc-editor.org/rfc/rfc9000}.

\end{thebibliography}

\end{document}
